\documentclass[12pt]{exam}
\usepackage{amsthm}
\usepackage{bm}
\usepackage{libertine}
\usepackage[utf8]{inputenc}
\usepackage[margin=0.5in]{geometry}
\usepackage{amsmath,amssymb}
\usepackage{multicol}
\usepackage[shortlabels]{enumitem}
\usepackage{siunitx}
\usepackage{physics}
\usepackage{booktabs}
\usepackage{graphicx}
\usepackage{pgfplots}
\usepackage{listings}
\usepackage{tikz}
\usepackage{tikz-3dplot}

\bmdefine{\ii}{i}
\bmdefine{\jj}{j}
\bmdefine{\kk}{k}
\newcommand{\te}{\theta}
\newcommand{\cC}{\mathcal{C}}
\newcommand{\word}[1]{\quad\text{#1}\quad}

\pgfplotsset{width=10cm,compat=1.9}
\usepgfplotslibrary{external}
\tikzexternalize

\newcommand{\class}{Math 261-001}
\newcommand{\examnum}{Quiz 1}
\newcommand{\examdate}{August 28}

\begin{document}
\pagestyle{plain}
\thispagestyle{empty}

\noindent
\textbf{\class}\\
\textbf{\examnum}, \textbf{\examdate} \\

\setlength{\tabcolsep}{3.5cm}
\renewcommand{\arraystretch}{1.5}
\setlength\extrarowheight{1cm}
\begin{tabular}{ |c|c| }
 \hline
 Name   & CSU ID \#  \\
 \hline
\end{tabular}

\vspace{10pt}

Be sure to read each question carefully. Answer all four problems. Each
problem is worth the same amount of points. Write your final answer
in the box provided. Only the answer written in the box will be graded. You
may use the surrounding space for scratch work. For a plane equation, any
nonzero scalar multiple of the same equation is equivalent.

\begin{enumerate}

\item Write $\vec{u}=(2,-1,3)$ in the form
$\vec{u}=a\ii+b\jj+c\kk$. \hfill [10 points]

\vfill
\begin{flushright}
\framebox(200,50){}
\end{flushright}

\item Let $\vec{u}=(1,2,-1)$ and
$\vec{v}=(3,-1,2)$. Compute
$\vec{u}\cdot\vec{v}$. \hfill [10 points]

\vfill
\begin{flushright}
\framebox(200,50){}
\end{flushright}

\newpage

\item Write a parametric equation $\vec{r}(t)$ for the line through
$P=(1,-1,2)$ and $Q=(3,1,1)$. \hfill [10 points]

\vfill
\begin{flushright}
\framebox(200,50){}
\end{flushright}

\item Let $P=(1,0,1)$, $Q=(2,1,1)$, and $R=(1,1,2)$.
\begin{enumerate}[(a)]
  \item Compute $\overrightarrow{PQ}\times\overrightarrow{PR}$.
  \hfill [5 points]

  \vfill
\begin{flushright}
\framebox(200,50){}
\end{flushright}
  \item Use the cross product from part (a) as a normal vector to write one
  equation of the plane through $P$, $Q$, and $R$. \hfill [5 points]
\vfill
\begin{flushright}
\framebox(200,50){}
\end{flushright}
\end{enumerate}


\end{enumerate}

\end{document}
